% part: intuitionistic-logic
% chapter: propositions-as-types
% section: types

\documentclass[../../../include/open-logic-section]{subfiles}

\begin{document}

\olfileid{int}{pty}{typ}

\olsection{Types}

A proof term such as $(MN)$ by itself does not uniquely determine the
proof it is a proof term of. That's because the free variables, i.e.,
the proof terms of axioms, do not specify the !!{formula}s in the
corresponding axioms. However, if we specify these in a
context~$\Gamma$, the proof \emph{is} uniquely determined.

\begin{defn}
A \emph{context} for a proof term~$N$ is a set $\Gamma$ such that for
every free variable of~$N$, $\Gamma$ contains exactly one pair $x : !A$.
\end{defn}

Note that the definition does not require that $\Gamma$ contains
$x:!A$ \emph{only} for free variables of~$N$. This, e.g., $\Gamma =
\{x:!A, y:!B, z:!C\}$ is a context for $\pair{x}{y}$.

\begin{defn}\ollabel{defn:type}
  Suppose $\Gamma$ is a context for~$N$. The \emph{type} of~$N$
  (relative to~$\Gamma$) is inductively defined as follows:
  \begin{enumerate}
  \item The type of $x$ is $!A$ iff $x:!A \in \Gamma$.
  \item If $N$ has type~$!B$ relative to $x:!A, \Gamma$, then the type of $\lambd[\typeof{x}{!A}][N]$ is $!A \lif !B$.
  \item If $N$ has type $!A \lif !B$ and $M$ has type~$!A$ then $(NM)$
  has type~$!B$
  \item If $N$ has type $!A$ and $M$ has type~$!B$, then $\pair{N}{M}$
  has type $!A \land !B$.
  \item If $N$ has type~$!A_1 \land !A_2$, then $\proj{i}{N}$ has type~$!A_i$.
  \item If $N$ has type $!A$, then
    $\inj[!B]{i}{!M}$ has type $!A \lor !B$ if $i=1$ and type $!B \lor
    !A$ if $i=2$.
  \item If $M$ has type $!A \lor !B$, and $N_1$ has type $!C$ relative
  to $x:!A, \Gamma$ and $N_2$ has type~$!C$ relative to $y:!B,
  \Gamma$, then $\dcase{M}{x_1}{N_1}{x_2}{N_2}$ has type~$!C$.
  \item If $N$ has type~$\lfalse$, then
    $\abort{!A}{N}$ has type~$!A$.
  \end{enumerate}
\end{defn}

\begin{prop}
Every proof term~$N$ has exactly one type relative to a
context~$\Gamma$ for~$N$.
\end{prop}

\begin{proof}
  By induction on~$N$.
\end{proof}

We write $\Gamma \Sequent N : !A$ if the type of $N$ relative to
$\Gamma$ is~$!A$. The clauses in \olref{defn:type} should seem familiar:
they are precisely the rules of \Log{tN2i} with one slight difference:
the context~$\Gamma$ is the same throughout. We can codify the
definition as a calculus \Log{tN3i} with the rules given in \olref{tab:tN2ip}.

\subfile{rules-tN3}

Proof terms are terms in a version of the $\lambd$-calculus, the
\emph{typed $\lambd$-calculus with products, sums, and empty type}. The
classical $\lambd$-calculus only has term built from variables,
\emph{$\lambd$-abstracts}~$\lambd[x][N]$ and application $(MN)$, and
it lacks the annotation by~$!A$ of the typed $\lambd$-abstract
$\lambd[\typeof{x}{!A}][N]$. The $\lambd$-calculus is a model of
computation which is Turing complete (i.e., every partial computable
function can be represented in it by a term). The version of the
$\lambd$-calculus given by proof terms and typing rules is a
restricted model of computation, but the basic ideas are the same. It
is restricted by the type assignment: only \emph{well-typed} terms,
i.e., proof terms with types (relative to a context~$\Gamma$) are
legal terms. For instance, $\lambd[x][(xx)]$ is a legal term in the
classical untyped $\lambd$-calculus, but
$\lambd[\typeof{x}{!A}][(xx)]$ is not well-typed for any $!A$. That's
because the typing rules require that in order for $(xx)$ to have a
type, the first~$x$ would have to have type $!A \lif !B$ and the
second $x$ type~$!A$, and that's impossible since $!A$ is a proper
sub-!!{formula} of~$!A \lif !B$.

We can think of types intuitively as the kind of thing named by a term
of that type. So a term of type $!A \land !B$ picks out a pair with
first component of type~$!A$ and second component of type~$!B$, i.e.,
!!a{element} of the product~$!A \times !B$. A term of type $!A \lif !B$
is a function from things of type $!A$ to things of type~$!B$. A term
of type $!A \lor !B$ is either something of type $!A$ or something of
type~$!B$, plus the information which one it is. This is like the
disjoint union of the two sets $!A$ and $!B$ defined as $\{\tuple{i,x}
: i = 1 \text{ or } 2, x \in !A \text{ if } i = 1, x \in !B \text{ if
} i=2\}$. The type $\lfalse$ is the empty set, i.e., if a term has
that type it can't name anything. Since natural deduction (including
\Log{tN3i}) is consistent, there is no !!{proof} of~$\lfalse$ unless
there are open assumptions, and so there is no closed proof term (term
without free variables) that has type~$\lfalse$.

The operators of the typed $\lambd$-calculus produce terms that
intuitively name the right sort of thing, provided the terms it is
applied to name certain things. E.g., ``if $N_1 : !A$ and $N_2 : !B$,
then $\pair{N_1}{N_2} : !A \land !B$'' says that if $N_1$ names
somethign of type~$!A$ and $N_2$ names something of type~$!B$, then
$\pair{N_1}{N_2}$ names something of type $!A \land !B$.

\end{document}
